% GENERATED FILE. Edit canonical lessons, then run npm run generate:books.
% canonical-source-hash: fnv1a64:b964546bfe19e66f
% canonical-lessons: HI-C499-kisor, HI-C499-balig, HI-C499-sararat, HI-C499-yauvan, HI-C499-bisurna

\chapter{Growing Up}
\label{ch:hi-growing-up}
\hlchaptermodality{499}

\begin{chapteropening}
\textbf{By the end of this chapter:} \emph{I can talk about growing up and coming of age.}

Complete the last lesson of chapter 499: I can talk about growing up and coming of age.
\end{chapteropening}

\section[\texorpdfstring{\hi{किशोर} (kiśor)}{किशोर (kiśor)}]{\hi{किशोर} (kiśor) — an adolescent; a teenager}

\label{lesson:HI-C499-kisor}



\begin{quote}
Before the new one: say the Hindi for to feed, then the Hindi for to be brought up.
\end{quote}

\subsection*{You'll want to know: \hi{किशोर}}
\textbf{\hi{किशोर}} — \emph{kiśor} — \textquotedblleft{}an adolescent; a teenager\textquotedblright{}.

\begin{quote}
\hi{किशोर} \hi{उम्र} \hi{में} \hi{बच्चे} \hi{बहुत} \hi{जल्दी} \hi{बदलते} \hi{हैं।}
\end{quote}

\begin{quote}
\emph{kiśor umr meṁ bacce bahut jaldī badalte haiṁ.}
\end{quote}

\begin{quote}
In their teens, children change very quickly.
\end{quote}

Masculine.

\begin{cousinweb}
\textbf{\hi{किशोर}} (\emph{kiśor}) is a tatsama, borrowed from Sanskrit \emph{kiśora}, first \textquotedblleft{}a colt\textquotedblright{} and then \textquotedblleft{}a youth, a lad\textquotedblright{} (Monier-Williams). The feminine is \emph{kiśorī}.
\end{cousinweb}

\begin{grammarlens}[title={What you've built}]
One more word for describing what you lived and felt.
\end{grammarlens}

\subsection*{Your turn}
\begin{itemize}
\raggedright
  \item \emph{Say it:} \emph{kiśor}
  \item \emph{Say it:} \emph{kiśor}, once more
  \item \emph{Say it:} \emph{kiśor umr meṁ bacce bahut jaldī badalte haiṁ.} — In their teens, children change very quickly.
  \item \emph{Recall:} say the Hindi for to feed, then the Hindi for to be brought up, then say \emph{kiśor} again
\end{itemize}

\begin{tcolorbox}[breakable,colback=teal!4,colframe=teal!35!black,title={Before you move on}]
What does \hi{किशोर} mean? (\textquotedblleft{}An adolescent; a teenager\textquotedblright{}.) Say it once more.
\end{tcolorbox}
\section[\texorpdfstring{\hi{बालिग़} (bāliġ)}{बालिग़ (bāliġ)}]{\hi{बालिग़} (bāliġ) — of age; adult}

\label{lesson:HI-C499-balig}



\begin{quote}
Before the new one: say the Hindi for to be brought up, then the Hindi for an adolescent.
\end{quote}

\subsection*{You'll want to know: \hi{बालिग़}}
\textbf{\hi{बालिग़}} — \emph{bāliġ} — \textquotedblleft{}of age; adult\textquotedblright{}.

\begin{quote}
\hi{अठारह} \hi{साल} \hi{का} \hi{होकर} \hi{अरुण} \hi{बालिग़} \hi{हो} \hi{गया।}
\end{quote}

\begin{quote}
\emph{aṭhārah sāl kā hokar Aruṇ bāliġ ho gayā.}
\end{quote}

\begin{quote}
At eighteen Arun came of age.
\end{quote}

\begin{cousinweb}
\textbf{\hi{बालिग़}} (\emph{bāliġ}) is usually traced to Arabic \emph{bāliġ}, \textquotedblleft{}having reached maturity\textquotedblright{}, which reached Hindi through Persian; it is the word of the law.
\end{cousinweb}

\begin{grammarlens}[title={What you've built}]
One more word for describing what you lived and felt.
\end{grammarlens}

\subsection*{Your turn}
\begin{itemize}
\raggedright
  \item \emph{Say it:} \emph{bāliġ}
  \item \emph{Say it:} \emph{bāliġ}, once more
  \item \emph{Say it:} \emph{aṭhārah sāl kā hokar Aruṇ bāliġ ho gayā.} — At eighteen Arun came of age.
  \item \emph{Recall:} say the Hindi for to be brought up, then the Hindi for an adolescent, then say \emph{bāliġ} again
\end{itemize}

\begin{tcolorbox}[breakable,colback=teal!4,colframe=teal!35!black,title={Before you move on}]
What does \hi{बालिग़} mean? (\textquotedblleft{}Of age; adult\textquotedblright{}.) Say it once more.
\end{tcolorbox}
\section[\texorpdfstring{\hi{शरारत} (śarārat)}{शरारत (śarārat)}]{\hi{शरारत} (śarārat) — mischief}

\label{lesson:HI-C499-sararat}



\begin{quote}
Before the new one: say the Hindi for an adolescent, then the Hindi for of age.
\end{quote}

\subsection*{You'll want to know: \hi{शरारत}}
\textbf{\hi{शरारत}} — \emph{śarārat} — \textquotedblleft{}mischief\textquotedblright{}.

\begin{quote}
\hi{बचपन} \hi{में} \hi{मीरा} \hi{बहुत} \hi{शरारत} \hi{करती} \hi{थी।}
\end{quote}

\begin{quote}
\emph{bacpan meṁ Mīrā bahut śarārat kartī thī.}
\end{quote}

\begin{quote}
As a child Mira got up to a lot of mischief.
\end{quote}

Feminine.

\begin{cousinweb}
\textbf{\hi{शरारत}} (\emph{śarārat}) is usually traced to Arabic \emph{sharāra}, \textquotedblleft{}wickedness\textquotedblright{}, which reached Hindi through Persian and softened to a child's mischief.
\end{cousinweb}

\begin{grammarlens}[title={What you've built}]
One more word for describing what you lived and felt.
\end{grammarlens}

\subsection*{Your turn}
\begin{itemize}
\raggedright
  \item \emph{Say it:} \emph{śarārat}
  \item \emph{Say it:} \emph{śarārat}, once more
  \item \emph{Say it:} \emph{bacpan meṁ Mīrā bahut śarārat kartī thī.} — As a child Mira got up to a lot of mischief.
  \item \emph{Recall:} say the Hindi for an adolescent, then the Hindi for of age, then say \emph{śarārat} again
\end{itemize}

\begin{tcolorbox}[breakable,colback=teal!4,colframe=teal!35!black,title={Before you move on}]
What does \hi{शरारत} mean? (\textquotedblleft{}Mischief\textquotedblright{}.) Say it once more.
\end{tcolorbox}
\section[\texorpdfstring{\hi{यौवन} (yauvan)}{यौवन (yauvan)}]{\hi{यौवन} (yauvan) — the prime of youth}

\label{lesson:HI-C499-yauvan}



\begin{quote}
Before the new one: say the Hindi for of age, then the Hindi for mischief.
\end{quote}

\subsection*{You'll want to know: \hi{यौवन}}
\textbf{\hi{यौवन}} — \emph{yauvan} — \textquotedblleft{}the prime of youth\textquotedblright{}.

\begin{quote}
\hi{दादी} \hi{कहती} \hi{हैं} \hi{कि} \hi{यौवन} \hi{जल्दी} \hi{बीत} \hi{जाता} \hi{है।}
\end{quote}

\begin{quote}
\emph{dādī kahtī haiṁ ki yauvan jaldī bīt jātā hai.}
\end{quote}

\begin{quote}
Grandmother says that youth passes quickly.
\end{quote}

Masculine.

\begin{cousinweb}
\textbf{\hi{यौवन}} (\emph{yauvan}) is a tatsama, borrowed from Sanskrit \emph{yauvana}, \textquotedblleft{}youth, youthfulness\textquotedblright{}, from \emph{yuvan}, \textquotedblleft{}young\textquotedblright{} (Monier-Williams); everyday speech says \emph{javānī}.
\end{cousinweb}

\begin{grammarlens}[title={What you've built}]
One more word for describing what you lived and felt.
\end{grammarlens}

\subsection*{Your turn}
\begin{itemize}
\raggedright
  \item \emph{Say it:} \emph{yauvan}
  \item \emph{Say it:} \emph{yauvan}, once more
  \item \emph{Say it:} \emph{dādī kahtī haiṁ ki yauvan jaldī bīt jātā hai.} — Grandmother says that youth passes quickly.
  \item \emph{Recall:} say the Hindi for of age, then the Hindi for mischief, then say \emph{yauvan} again
\end{itemize}

\begin{tcolorbox}[breakable,colback=teal!4,colframe=teal!35!black,title={Before you move on}]
What does \hi{यौवन} mean? (\textquotedblleft{}The prime of youth\textquotedblright{}.) Say it once more.
\end{tcolorbox}
\section[\texorpdfstring{\hi{बिसूरना} (bisūrnā)}{बिसूरना (bisūrnā)}]{\hi{बिसूरना} (bisūrnā) — to whimper; to sulk tearfully}

\label{lesson:HI-C499-bisurna}



\begin{quote}
Before the new one: say the Hindi for mischief, then the Hindi for the prime of youth.
\end{quote}

\subsection*{You'll want to know: \hi{बिसूरना}}
\textbf{\hi{बिसूरना}} — \emph{bisūrnā} — \textquotedblleft{}to whimper; to sulk tearfully\textquotedblright{}.

\begin{quote}
\hi{छोटा} \hi{भाई} \hi{मिठाई} \hi{के} \hi{लिए} \hi{बिसूर} \hi{रहा} \hi{था।}
\end{quote}

\begin{quote}
\emph{choṭā bhāī miṭhāī ke liye bisūr rahā thā.}
\end{quote}

\begin{quote}
My little brother was whimpering for a sweet.
\end{quote}

Past (he …): \hi{वह} \hi{दिन} \hi{भर} \hi{बिसूरा}, \emph{vah din bhar bisūrā}.

\begin{cousinweb}
\textbf{\hi{बिसूरना}} (\emph{bisūrnā}) is a tadbhava: CDIAL lists Hindi \emph{bisūrnā}, \textquotedblleft{}to put on a crying face\textquotedblright{}, under \emph{*viśūratē}, a form it reconstructs (CDIAL 11941).
\end{cousinweb}

\begin{grammarlens}[title={What you've built}]
One more word for describing what you lived and felt.
\end{grammarlens}

\subsection*{Your turn}
\begin{itemize}
\raggedright
  \item \emph{Say it:} \emph{bisūrnā}
  \item \emph{Say it:} \emph{bisūrnā}, once more
  \item \emph{Say it:} \emph{choṭā bhāī miṭhāī ke liye bisūr rahā thā.} — My little brother was whimpering for a sweet.
  \item \emph{Recall:} say the Hindi for mischief, then the Hindi for the prime of youth, then say \emph{bisūrnā} again
\end{itemize}

\begin{tcolorbox}[breakable,colback=teal!4,colframe=teal!35!black,title={Before you move on}]
What does \hi{बिसूरना} mean? (\textquotedblleft{}To whimper; to sulk tearfully\textquotedblright{}.) Say it once more.
\end{tcolorbox}
